\section*{Notation}
\addcontentsline{toc}{section}{Notation}

\emph{Typesetter's note:} see \verb|NOTATION.md| for the specific notation
upgrades I implemented. Some include switching to Dirac notation, switching
$a(\mathbf{p}) \to a_{\mathbf{p}}$, and getting rid of the terrible-to-look-at
\verb|mathscr| letters, e.g. $\mathscr{J} \to \mathcal{J}$.

Latin indices $i,j,k,\ldots$ generally run over the three spatial coordinates
$1,2,3$. Greek indices $\mu,\nu,\ldots$ run over spacetime coordinates in the
order $0,1,2,3$, with $x^0$ the time coordinate. Repeated indices are summed
unless otherwise indicated.

The spacetime metric is mostly plus,
\[
  \eta_{\mu\nu}=\operatorname{diag}(-1,+1,+1,+1),
\]
so an on-shell four-momentum satisfies $p^2=-m^2$, and spacelike separations
satisfy $(x-x')^2\geq 0$. The d'Alembertian is
\[
  \Box=\partial_\mu\partial^\mu=\nabla^2-\partial_t^2.
\]
The Levi--Civita tensor is totally antisymmetric, with
$\epsilon^{0123}=+1$.

Spatial three-vectors are bold, and a hat denotes a unit vector:
$\widehat{\mathbf v}=\mathbf v/\lvert\mathbf v\rvert$. A dot denotes a time
derivative.

Dirac matrices obey
\[
  \{\gamma^\mu,\gamma^\nu\}=2\eta^{\mu\nu},
  \qquad
  \gamma^5=i\gamma^0\gamma^1\gamma^2\gamma^3.
\]
With this metric convention, $\beta=i\gamma^0$ and the Dirac adjoint is
$\bar u=u^\dagger\beta$. The step function $\theta(s)$ equals $1$ for $s>0$
and $0$ for $s<0$.

Complex conjugation, transpose, and Hermitian conjugation are written $A^*$,
$A^{\mathsf T}$, and $A^\dagger$, respectively. The identity is $\mathbf{1}$.
The abbreviations $+\mathrm{H.c.}$ and $+\mathrm{c.c.}$ indicate addition of the
Hermitian conjugate and complex conjugate of the preceding terms.

States use Dirac notation: $\ket{p,\sigma}$ for one-particle states,
$\ket{\Omega}$ for the vacuum, and $\braket{\phi}{\psi}$ for inner products.
Asymptotic states are labeled $\ket{\alpha}_{\mathrm{in}}$ and
$\ket{\beta}_{\mathrm{out}}$. Positive- and negative-frequency field pieces
are written $\psi^{(+)}$ and $\psi^{(-)}$, avoiding the original overloading of
$+$ and $-$ for both field frequencies and asymptotic states.

The interaction Hamiltonian is denoted
\[
  H_I(t)=\int d^3x\,\mathcal H_I(x),
\]
in place of the original bare $V$ notation. Other operators are generally left
unhatted when the context is unambiguous.

Except in Chapter 1, $\hbar=c=1$. Throughout, $-e$ is the rationalized charge
of the electron, so $\alpha=e^2/(4\pi)\simeq 1/137$. Parentheses at the end of
quoted numerical data give the uncertainty in the final digits.
